% GENERATED FILE. Edit canonical lessons, then run npm run generate:books.
% canonical-source-hash: fnv1a64:2dc60e0c9e383a70
% canonical-lessons: AR-C310-istaghatha, AR-C310-naja, AR-C310-saha, AR-C310-nada, AR-C310-munqidh

\chapter{Calling for help}
\label{ch:ar-calling-for-help}
\hlchaptermodality{310}

\begin{chapteropening}
\textbf{By the end of this chapter:} \emph{I can say to call for help, to survive, to shout, to call out and a rescuer.}

Complete the last lesson of chapter 310: I can say to call for help, to survive, to shout, to call out and a rescuer.
\end{chapteropening}

\section[\texorpdfstring{\ar{استغاث} (istaghātha)}{استغاث (istaghātha)}]{\ar{استغاث} (istaghātha) — to call for help}

\label{lesson:AR-C310-istaghatha}



\begin{quote}
Before the new one: say the Arabic for a throne, then the Arabic for to crown.
\end{quote}

\subsection*{You'll want to know: \ar{استغاث}}
\textbf{\ar{استغاث}} — \emph{istaghātha} — \textquotedblleft{}to call for help\textquotedblright{}.

\begin{quote}
\ar{استغاثت} \ar{المرأة} \ar{من} \ar{نافذة} \ar{البيت}.
\end{quote}

\begin{quote}
\emph{istaghāthat al-marʾa min nāfidhat al-bayt.}
\end{quote}

\begin{quote}
The woman called for help from the window of the house.
\end{quote}

Present: \ar{يستغيث}, \emph{yastaghīth}.

\begin{cousinweb}
\textbf{\ar{استغاث}} (\emph{istaghātha}) is built on the native Arabic root \emph{gh-w-th} (Lane), where \emph{ghawth} is help given to one in distress. It is Form X, \textquotedblleft{}to ask for help\textquotedblright{}.
\end{cousinweb}

\begin{grammarlens}[title={What you've built}]
One more word for telling what happened.
\end{grammarlens}

\subsection*{Your turn}
\begin{itemize}
\raggedright
  \item \emph{Say it:} \emph{istaghātha}
  \item \emph{Say it:} \emph{istaghātha}, once more
  \item \emph{Say it:} \emph{istaghāthat al-marʾa min nāfidhat al-bayt.} — The woman called for help from the window of the house.
  \item \emph{Recall:} say the Arabic for a throne, then the Arabic for to crown, then say \emph{istaghātha} again
\end{itemize}

\begin{tcolorbox}[breakable,colback=teal!4,colframe=teal!35!black,title={Before you move on}]
What does \ar{استغاث} mean? (\textquotedblleft{}To call for help\textquotedblright{}.) Say it once more.
\end{tcolorbox}
\section[\texorpdfstring{\ar{نجا} (najā)}{نجا (najā)}]{\ar{نجا} (najā) — to survive, to escape}

\label{lesson:AR-C310-naja}



\begin{quote}
Before the new one: say the Arabic for to crown, then the Arabic for to call for help.
\end{quote}

\subsection*{You'll want to know: \ar{نجا}}
\textbf{\ar{نجا}} — \emph{najā} — \textquotedblleft{}to survive, to escape\textquotedblright{}.

\begin{quote}
\ar{نجا} \ar{كل} \ar{الركاب} \ar{من} \ar{الحادث}.
\end{quote}

\begin{quote}
\emph{najā kull ar-rukkāb min al-ḥādith.}
\end{quote}

\begin{quote}
All the passengers survived the accident.
\end{quote}

Present: \ar{ينجو}, \emph{yanjū}.

It takes \ar{من}: \ar{نجا} \ar{من} \ar{الموت}.

\begin{cousinweb}
\textbf{\ar{نجا}} (\emph{najā}) is built on the native Arabic root \emph{n-j-w} (Lane), \textquotedblleft{}to escape, to be safe\textquotedblright{}; \emph{najāt}, rescue, is its verbal noun.
\end{cousinweb}

\begin{grammarlens}[title={What you've built}]
One more word for telling what happened.
\end{grammarlens}

\subsection*{Your turn}
\begin{itemize}
\raggedright
  \item \emph{Say it:} \emph{najā}
  \item \emph{Say it:} \emph{najā}, once more
  \item \emph{Say it:} \emph{najā kull ar-rukkāb min al-ḥādith.} — All the passengers survived the accident.
  \item \emph{Recall:} say the Arabic for to crown, then the Arabic for to call for help, then say \emph{najā} again
\end{itemize}

\begin{tcolorbox}[breakable,colback=teal!4,colframe=teal!35!black,title={Before you move on}]
What does \ar{نجا} mean? (\textquotedblleft{}To survive, to escape\textquotedblright{}.) Say it once more.
\end{tcolorbox}
\section[\texorpdfstring{\ar{صاح} (ṣāḥa)}{صاح (ṣāḥa)}]{\ar{صاح} (ṣāḥa) — to cry out, to shout}

\label{lesson:AR-C310-saha}



\begin{quote}
Before the new one: say the Arabic for to call for help, then the Arabic for to survive.
\end{quote}

\subsection*{You'll want to know: \ar{صاح}}
\textbf{\ar{صاح}} — \emph{ṣāḥa} — \textquotedblleft{}to cry out, to shout\textquotedblright{}.

\begin{quote}
\ar{صاح} \ar{الولد} \ar{عندما} \ar{رأى} \ar{الكلب}.
\end{quote}

\begin{quote}
\emph{ṣāḥa al-walad ʿindamā raʾā al-kalb.}
\end{quote}

\begin{quote}
The boy shouted when he saw the dog.
\end{quote}

Present: \ar{يصيح}, \emph{yaṣīḥ}.

A cockerel \ar{صاح} too: it crowed.

\begin{cousinweb}
\textbf{\ar{صاح}} (\emph{ṣāḥa}) is built on the native Arabic root \emph{ṣ-y-ḥ} (Lane), \textquotedblleft{}to cry out, to raise the voice\textquotedblright{}; \emph{ṣayḥa} is a shout.
\end{cousinweb}

\begin{grammarlens}[title={What you've built}]
One more word for telling what happened.
\end{grammarlens}

\subsection*{Your turn}
\begin{itemize}
\raggedright
  \item \emph{Say it:} \emph{ṣāḥa}
  \item \emph{Say it:} \emph{ṣāḥa}, once more
  \item \emph{Say it:} \emph{ṣāḥa al-walad ʿindamā raʾā al-kalb.} — The boy shouted when he saw the dog.
  \item \emph{Recall:} say the Arabic for to call for help, then the Arabic for to survive, then say \emph{ṣāḥa} again
\end{itemize}

\begin{tcolorbox}[breakable,colback=teal!4,colframe=teal!35!black,title={Before you move on}]
What does \ar{صاح} mean? (\textquotedblleft{}To cry out, to shout\textquotedblright{}.) Say it once more.
\end{tcolorbox}
\section[\texorpdfstring{\ar{نادى} (nādā)}{نادى (nādā)}]{\ar{نادى} (nādā) — to call out to}

\label{lesson:AR-C310-nada}



\begin{quote}
Before the new one: say the Arabic for to survive, then the Arabic for to cry out.
\end{quote}

\subsection*{You'll want to know: \ar{نادى}}
\textbf{\ar{نادى}} — \emph{nādā} — \textquotedblleft{}to call out to\textquotedblright{}.

\begin{quote}
\ar{نادت} \ar{الأم} \ar{ابنها} \ar{من} \ar{الشرفة}.
\end{quote}

\begin{quote}
\emph{nādat al-umm ibnahā min ash-shurfa.}
\end{quote}

\begin{quote}
The mother called to her son from the balcony.
\end{quote}

Present: \ar{ينادي}, \emph{yunādī}.

Not the taught \ar{نادي}, a club: the final letter tells them apart.

\begin{cousinweb}
\textbf{\ar{نادى}} (\emph{nādā}) is built on the native Arabic root \emph{n-d-w} (Lane), \textquotedblleft{}to call; to meet together\textquotedblright{}, the root of the taught \emph{nādī}, a club. It is Form III, \textquotedblleft{}to call to someone\textquotedblright{}.
\end{cousinweb}

\begin{grammarlens}[title={What you've built}]
One more word for telling what happened.
\end{grammarlens}

\subsection*{Your turn}
\begin{itemize}
\raggedright
  \item \emph{Say it:} \emph{nādā}
  \item \emph{Say it:} \emph{nādā}, once more
  \item \emph{Say it:} \emph{nādat al-umm ibnahā min ash-shurfa.} — The mother called to her son from the balcony.
  \item \emph{Recall:} say the Arabic for to survive, then the Arabic for to cry out, then say \emph{nādā} again
\end{itemize}

\begin{tcolorbox}[breakable,colback=teal!4,colframe=teal!35!black,title={Before you move on}]
What does \ar{نادى} mean? (\textquotedblleft{}To call out to\textquotedblright{}.) Say it once more.
\end{tcolorbox}
\section[\texorpdfstring{\ar{منقذ} (munqidh)}{منقذ (munqidh)}]{\ar{منقذ} (munqidh) — a rescuer; a lifeguard}

\label{lesson:AR-C310-munqidh}



\begin{quote}
Before the new one: say the Arabic for to cry out, then the Arabic for to call out to.
\end{quote}

\subsection*{You'll want to know: \ar{منقذ}}
\textbf{\ar{منقذ}} — \emph{munqidh} — \textquotedblleft{}a rescuer; a lifeguard\textquotedblright{}.

\begin{quote}
\ar{سبح} \ar{المنقذ} \ar{نحو} \ar{الطفل} \ar{بسرعة}.
\end{quote}

\begin{quote}
\emph{sabaḥa al-munqidh naḥwa aṭ-ṭifl bi-surʿa.}
\end{quote}

\begin{quote}
The lifeguard swam quickly towards the child.
\end{quote}

\begin{cousinweb}
\textbf{\ar{منقذ}} (\emph{munqidh}) is built on the native Arabic root \emph{n-q-dh} (Lane), the root of the taught \emph{anqadha}, to rescue. \emph{Munqidh} is its Form IV participle, \textquotedblleft{}one who rescues\textquotedblright{}.
\end{cousinweb}

\begin{grammarlens}[title={What you've built}]
One more word for telling what happened.
\end{grammarlens}

\subsection*{Your turn}
\begin{itemize}
\raggedright
  \item \emph{Say it:} \emph{munqidh}
  \item \emph{Say it:} \emph{munqidh}, once more
  \item \emph{Say it:} \emph{sabaḥa al-munqidh naḥwa aṭ-ṭifl bi-surʿa.} — The lifeguard swam quickly towards the child.
  \item \emph{Recall:} say the Arabic for to cry out, then the Arabic for to call out to, then say \emph{munqidh} again
\end{itemize}

\begin{tcolorbox}[breakable,colback=teal!4,colframe=teal!35!black,title={Before you move on}]
What does \ar{منقذ} mean? (\textquotedblleft{}A rescuer; a lifeguard\textquotedblright{}.) Say it once more.
\end{tcolorbox}
